\subsection{Höhere Ableitungen}

Wir definieren zunächst die Funktionalklassen, die angeben, wie oft eine Funktion partiell differenzierbar ist und ob ihre Ableitungen stetig sind:


\[
    C^0(D) \equiv C^0(D,  \mathbb{R}) := \{\,f: D\to\mathbb{R}\mid f \text{ ist stetig}\},
\]




\[
C^1(D)=\{\,f: D\to\mathbb{R}\mid f \text{ ist partiell differenzierbar und alle partiellen Ableitungen sind stetig}\},
\]




\[
C^k(D)=\Bigl\{\,f: D\to\mathbb{R}\;\Big|\; f \text{ ist } k\text{-mal partiell differenzierbar und alle Ableitungen bis zur Ordnung } k \text{ sind stetig}\Bigr\}.
\]


Insbesondere bezeichnet \(C^\infty(D)\) die Menge der unendlich oft differenzierbaren Funktionen:
\[
    C^\infty(D) \equiv C^\infty(C,\mathbb{R}) := \bigcap_{k=1}^\infty C^k(D)
\]

\begin{bemerkungbox}
Wir verwenden häufig die Notation \(D^k f\) oder \(\partial^{(k)} f\) für die \(k\)-fachen partiellen Ableitungen.
\end{bemerkungbox}

\begin{satz}{}
\textbf{(Schwarzsche Ableitungsregel)}\\
Sei \(f\in C^2(D)\) für eine offene Menge \(D\subset \mathbb{R}^n\). Dann gilt für alle Punkte \(x\in D\) und für alle Indizes \(i,j\in\{1,\dots,n\}\) 


\[
\frac{\partial^2 f}{\partial x_i\partial x_j}(x)=\frac{\partial^2 f}{\partial x_j\partial x_i}(x).
\]


\end{satz}

\begin{beweisbox}
Der Beweis erfolgt zunächst im Fall \(n=2\) durch Anwendung des Mittelwertsatzes auf die Differenzquotienten und lässt sich anschließend induktiv auf höhere Dimensionen verallgemeinern. (Details entnehmen Sie bitte der weiterführenden Literatur.)
\end{beweisbox}

\begin{bemerkungbox}
Ist \(f\) lediglich zweifach partiell differenzierbar, so kann es vorkommen, dass die gemischten Ableitungen ohne zusätzliche Stetigkeitsannahmen nicht übereinstimmen. Die Stetigkeit der zweiten Ableitungen – also \(f\in C^2(D)\) – ist daher wesentlich für die Gültigkeit der Schwarzschen Regel.
\end{bemerkungbox}

\begin{korollar}
Sei \(f\in C^2(D)\). Dann ist die Hessesche Matrix


\[
H_f(x)=\Bigl( \frac{\partial^2 f}{\partial x_i \partial x_j}(x)\Bigr)_{i,j=1}^n
\]


für alle \(x\in D\) symmetrisch.
\end{korollar}

\begin{bemerkungbox}
Die Konzepte der totalen Ableitung und der Hesseschen Matrix lassen sich auch auf Funktionen zwischen Banachräumen verallgemeinern. In diesem Rahmen erhält man analoge Resultate bezüglich der Stetigkeit der Ableitungen.
\end{bemerkungbox}
